% intro.tex — Introduction body, \input from septagraph.tex.
% Draft for emgor to re-voice. Beats, in order:
%   1 the concrete problem (one grid, seven robots)
%   2 why note/channel vocabularies fail
%   3 MIDI's actual lesson (underspecification)
%   4 the generative-audio control gap
%   5 thesis (separation; where / how tightly / when)
%   6 contributions + the interchangeability matrix
% ACTION(emgor): confirm authorship/acknowledgment with Kapur and Colton
%   before submission (the Machine Lab is named here and in Sec. origin).

% --- 1. The problem -------------------------------------------------
The Machine Lab at CalArts is a room of seven robotic instruments ---
BreakBot, GanaPati, Tammy, Marimba, Trimpbeat, Modulettes, and
RattleTron --- some tonal, some purely percussive, each bolted to its
own spot in the room. Composers drive them today as the ensemble has
always been driven: every instrument is a separate endpoint, reached
by its own OSC address behind a MIDI channel and note number that
select a beater and a velocity. Each machine has mechanical limits the
sender must already know --- a maximum re-strike rate, a force range,
and, on GanaPati, a thermal cutout: driven too fast it overheats and
shuts down, so faster play must arrive at lower velocity or not at
all. Nothing in the message vocabulary carries any of this.

We wanted to play the room as \emph{one} instrument from one $8\times8$
grid: move a hand across it and have the sound move from drum to
marimba, dry to noisy, left to right in the room, without ever
deciding which robot to address.

% --- 2. Why notes and channels fail ---------------------------------
Note and channel vocabularies cannot express that request. In MIDI, a
robot is a channel and a beater is a note number; ``a bit noisier and
a little further left'' is not a message that exists. There is no
continuum between a drum and a marimba because there is no axis along
which they differ --- they are two channels. The same wall stands
between two software synthesizers, or between a synthesizer and a
robot: as long as the shared language names \emph{sources}, a
controller can only choose one.

% --- 3. MIDI's actual lesson ----------------------------------------
It is tempting to conclude the fix is a richer description of sound.
MIDI's history argues the opposite. MIDI became universal because it
said almost nothing --- which key, how hard, when --- and left the
instrument to own everything else. That underspecification is why a
1983 keyboard still drives a 2026 synthesizer. A control standard
survives by naming the smallest thing both sides can agree on. MIDI's
problem is not that it is thin; it is that the thin thing it names ---
a key on a keyboard --- stopped being the right primitive once the
receiver was no longer a keyboard instrument.

% --- 4. The generative-audio control gap ----------------------------
The gap is now acute. Neural and generative sound engines have no
control vocabulary beyond whatever each model happens to expose: a
text prompt, or a latent vector meaningless to any other model.
Neither is portable; neither is something a hand can play in real
time. What \emph{is} portable, across a neural engine, a deterministic
synthesizer, and a solenoid-struck drum, is a perceptual description
of the result --- how bright, how noisy, how fast it flutters, where
in the room it sits. These are the quantities a listener hears and a
performer aims for, and any engine can be asked to reach a point on
them in its own way.

% --- 5. Thesis ------------------------------------------------------
Septagraph is built on one separation: the side that \emph{predicts
the human} is kept apart from the side that \emph{renders the sound},
meeting only in a shared perceptual coordinate space expressed as an
OSC sublanguage. A controller says \emph{where} in that space a sound
should land, \emph{how tightly} ($\sigma$, the renderer's declared
freedom to deviate), and \emph{when} --- discrete events on a fast
plane no model ever touches, and slower reconfiguration delivered only
at gesture boundaries, so whatever intelligence sits behind the
controller can be arbitrarily slow without its latency ever being
heard. A renderer announces what it can reach: a continuous region for
a synthesizer, or a measured \emph{catalog} of discrete positions for a
physical ensemble, each carrying the mechanical constraints the robots
forced us to model. Neither side knows anything else about the other.

% --- 6. Contributions and the matrix --------------------------------
This paper contributes: (i) the coordinate space and its wire grammar
(Secs.~\ref{sec:space}--\ref{sec:window}); (ii) a zero-install browser
controller/renderer pair on WebChucK a reader can run without
downloading anything; (iii) a Machine Lab renderer that speaks the
protocol into the ensemble's existing OSC server from a catalog of
measured robot positions; (iv) the neuralGrid, a physical grid, ribbon,
and lidar controller; (v) a latency evaluation of the fast plane; and
(vi) a listening experiment testing whether the axes are independently
steerable and whether the same coordinate stream is recognizable across
renderers. The central evidence is an interchangeability matrix ---
controllers against renderers, every cell the same protocol stream and
a working performance (Sec.~\ref{sec:eval}). A neural renderer
conditioned on the same axes is in progress with collaborators and is
discussed as future work. Before the technical sections, we describe
where this came from, because the robots shaped most of what follows.
